\section*{Introduction}

Cette note examine quatre réformes avec un seul critère : à recette à peu près constante, quel système
fiscal gêne le moins la croissance, c'est-à-dire décourage le moins de travailler, d'investir, de déménager
ou de déclarer ses revenus ? La redistribution n'est pas l'objectif, mais deux contraintes s'imposent. La
somme des réformes ne doit pas alourdir nettement la facture des plus grandes fortunes : leur départ coûterait
plus qu'il ne rapporterait. Et le système doit être, à la marge, plus favorable aux familles avec enfants.
Il s'agit de taxer mieux, pas davantage.

La méthode est la même pour les quatre impôts. On part d'une référence simple : l'assiette la plus large,
le même taux pour tout le monde, sans niche. Puis on fait apparaître chacun des endroits où un choix doit
être fait (une franchise pour les revenus d'activité, un traitement des familles, la vitesse de la
déduction de l'investissement, le sort des plus-values logées dans une holding) et on chiffre chaque
choix de la même façon : ce qu'il coûte en recette, ce qu'il coûte ou rapporte en perte sèche, et qui
paie. La note ne tranche pas ces arbitrages à la place du lecteur ; elle dit ce que chacun coûte.

\begin{encadre}{Comment la perte sèche est calculée}
\textbf{Ce qu'elle mesure.} Un impôt fait renoncer à des choses qui auraient valu plus qu'elles ne coûtaient :
une heure de travail, une machine, une vente de logement. La perte sèche est la valeur de ces renoncements,
au-delà de l'impôt lui-même. Elle se mesure par la réaction des quantités taxées à l'impôt.

\textbf{La formule.} La perte sèche ordinaire, ½·ε·τ²/(1$-$τ)·base, suppose une seule marge et aucun autre impôt
dessus. On utilise sa généralisation (Harberger ; Hines, 1999 ; Chetty, 2009) : la variation de perte sèche
est la somme, sur chaque marge touchée, du coin fiscal total qui frappe cette marge multiplié par la
variation de la quantité taxée. Un euro de salaire perdu coûte à la collectivité tout ce qu'il aurait payé :
impôt sur le revenu et CSG, mais aussi cotisations non contributives et prestations perdues. Trois
conséquences comptent :
\begin{itemize}\setlength{\itemsep}{0pt}
\item \emph{le coin total} : relever l'impôt là où les cotisations et le retrait des prestations prélèvent
déjà 60~\% d'un euro de salaire en plus (autour de 1 à 1,5 Smic) coûte beaucoup plus que le calcul sur le
seul impôt ne le dit ;
\item \emph{les déplacements entre impôts} : un revenu qui quitte le barème pour l'impôt sur les sociétés ou le
prélèvement forfaitaire n'est perdu que pour l'écart de taux ; une réaction de calendrier (distribuer un
dividende un an plus tard) ne coûte presque rien ;
\item \emph{la décision de travailler} : on compte aussi la réaction de ceux qui hésitent à prendre un emploi,
qui dépend de ce que l'impôt prélève sur l'ensemble du premier salaire.
\end{itemize}

\textbf{Comparer des variantes.} Quand deux variantes ne lèvent pas la même recette, la différence doit être
levée ailleurs. On la valorise au « prix d'un milliard » levé par le barème fusionné : 0,16 à 0,60~Md€ de perte
sèche par milliard selon le coin retenu, 0,38 au centre (méthode et paramètres : \texttt{revue/methode\_perte\_seche.md}).

\textbf{Ce qui reste fragile.} Les élasticités, et surtout deux paramètres sans estimation française : la part
de la réaction des hauts revenus qui n'est qu'un déplacement vers l'impôt sur les sociétés, et la part de
l'ajustement des distributions de dividendes qui n'est qu'un report. Chaque résultat est donné avec sa
fourchette.
\end{encadre}
